\section{Phase-count completeness from totality}\label{sec:completeness}
The halting reduction preserves all-law ergodicity because its controller
forgets every start. For total-count completeness, stationary
common-noise couplings suffice. This section uses the input categories
of Section~\ref{sec:framework} with the rate replaced by
\begin{equation}\label{eq:totality-rate}
 \eps_1=160000^{-101}.
\end{equation}

\subsection{The source condition}
\emph{Condition \(C_T\)} is the validity of the repaired
T\"orm\"a-based totality source with the following interface, using
\cite[Lemma~25, Corollary~26 and Lemmas~27--29]{torma} as published.
There is a total computable family of one-dimensional radius-one
CA \(U_e\), with finite alphabet \(\mathsf D_e\), blank \(0\), and
local table \(u_e:\mathsf D_e^3\to\mathsf D_e\), such that:
\begin{enumerate}
\item \(u_e(0,0,0)=0\). If the partial computable function
\(\varphi_e\) is total, every fully \(\Z^2\)-TI law on the
two-sided input-erasure trajectory SFT is all blank.
\item If \(\varphi_e\) is not total, the ordinary two-sided
trajectory SFT has a fully \(\Z^2\)-TI law with positive
nonblank activity.
\end{enumerate}
Here input erasure means replacing selected parent symbols by \(0\)
from one common preceding row. The repaired source reserves zero
for each bookkeeping field's blank value, normalizes invalid own
bookkeeping to blank, leaves simulated data unrestricted, and
increments age only while live; field widths are fixed before
self-reference. Quiescence of the universal layer is part of this
source premise. The interface is stronger than deterministic
unique ergodicity. The cited published lemmas supply the macrocell,
overlap, rigidity, connecting and live-coverage ingredients; their
internal proofs are used as published.

\subsection{A controller with diagonal stationary couplings}
A canonical common-noise pair uses the same replacement flag and
replacement letter in both replicas at each site and time. Say a
controller has the \emph{diagonal-coupling property} if every TI
stationary law of this pair kernel is supported on equal replicas.

\begin{lemma}\label{lem:totality-controller}
Under \(C_T\), one can compute from \(e\) a deterministic rule
\(G_e\) on \(\Z^6\), of maximum-norm radius at most \(11\), followed
by uniform replacement at \(\eps_1\), such that totality gives the
diagonal-coupling property and nontotality gives two distinct TI
stationary laws.
\end{lemma}
\begin{proof}
\emph{Certificates and quorum.}
Put
\[
 \mathsf L_e=\{0\}\sqcup
 \{(a,w):a\ne0,\ w\in\mathsf D_e^3,\ u_e(w)=a\},
 \qquad \operatorname{sym}(0)=0,\quad
 \operatorname{sym}(a,w)=a.
\]
On \(u=(i,s)\in\Z^2\), let \(H_{0,e}\) retain a nonblank
label \((a,w)\) exactly when every nonzero \(w_j\) agrees with
the symbol at \((i+j,s-1)\), \(j=-1,0,1\); otherwise delete it.
Blank centers remain blank. A fixed field projects to an
input-erasure trajectory, using its stored witness at a nonblank
center and erasing all parents at a blank one. Conversely, an
ordinary trajectory lifts by storing its actual predecessor word.

Physical sites are \((u,P)\in\Z^2\times\Z^4\).
Let \(H_e\) require, in addition, equality of the full label with
all eight transverse unit neighbors. Its stencil is
\[
 \mathcal N_*=
 \bigl(\{0,(-1,-1),(0,-1),(1,-1)\}\times\{0\}\bigr)
 \cup\{(0,\pm e_j):1\le j\le4\}.
\]
Use providers at \(c_0=0\), \(c_j=10e_j\), \(1\le j\le4\).
The rule \(G_e\) returns a nonblank label if four of the five
provider values of \(H_e\) equal it, and returns \(0\) otherwise.
The \(60\) read offsets are distinct and have radius at most \(11\).
Since \(|\mathsf L_e|\le1+|\mathsf D_e|^3\), its finite table is
computable from the source table.

\emph{Noise marks.}
Write \(\mathsf{Flag}_t(v)\) for the iid replacement flags. Mark a site if
it is flagged or if two distinct providers each have a marked
\(\mathcal N_*\)-parent. Empty starts receding to the past give
the least two-sided mark field \(E^*\), a stationary causal
factor of the flags. Every mark has a finite support graph
ending at distinct physical flags.

Here are the flag estimates used below. The backward graph has
undirected degree at most \(120\); use \(d_*=200\). The five
potentials
\[
 \psi_j(P,t)=-P_j-2t\quad(1\le j\le4),\qquad
 \psi_0(P,t)=\sum_{j=1}^4P_j+8t
\]
sum to zero. Provider \(j\) is excluded for \(\psi_j\),
\(0\le j\le4\). Along a backward provider step the gain is at least
\(1\) except at the excluded provider, and every gain
is at least \(-9\). Two distinct providers therefore offer a
good step for each potential.

To count a finite support graph while reusing terminals correctly,
group frontier vertices by overlap of their descendant-terminal
sets. Exposing the two supports of one vertex can split its
component into at most two and cannot merge different components.
Start with one pole at the root, and keep one terminal endpoint
of each potential in each component.
At a split add a pole, sending its arm of each potential to the
daughter lacking that endpoint. Only that first edge may be bad.
In the original root's arm component, \(n_{\mathrm f}\) distinct
terminal flags give \(n_{\mathrm f}\) poles and at most
\(5(n_{\mathrm f}-1)\) bad edges. For each potential the arms
biject these poles with the terminal flags. The sum of all arm
gains is zero, since the potentials balance at each pole and
terminal. Telescoping bounds the total number of arm edges by
\(50(n_{\mathrm f}-1)\): replacing a gain of at least \(1\) by
one of at least \(-9\) loses at most \(10\).
Physical projection cannot increase
edge count or time span.

A rooted connected graph with \(a\) edges has at most
\(d_*^{2a}\) doubled-walk encodings and \(2^{a+1}\) choices of
terminal subset. Distinct flags cost \(\eps_1^{n_{\mathrm f}}\), even when
several arms reach the same flag. Thus, with
\[
 \lambda=2d_*^2\eps_1^{1/50}
       =\tfrac12\,160000^{-51/50}<\tfrac12,
\]
the mark density satisfies
\begin{equation}\label{eq:mark-density}
 p_*:=\Pr(E^*(0)=1)
 \le\eps_1+2\eps_1\sum_{a\ge1}\lambda^a
 <3\eps_1<\tfrac12.
\end{equation}
For a marked nonflagged backward path of length \(n\), take complete
support graphs in successive disjoint time slabs, restarting below
the earliest terminal of each and keeping the last graph complete.
Their distinct \(n_{\mathrm f}\) flags cover \(n\le50n_{\mathrm f}\)
path steps. The connected union has \(n\le a\le100n_{\mathrm f}\)
edges. The same counting gives
\begin{equation}\label{eq:marked-path}
 \Pr(\text{such a path from a fixed root})
 \le \frac{2\gamma_*^n}{1-\gamma_*},\qquad
 \gamma_*=2d_*^2\eps_1^{1/100}
         =\tfrac12\,160000^{-1/100}<\tfrac12.
\end{equation}
Consequently infinite marked nonflagged backward paths have
probability zero, simultaneously at all countably many roots.

\emph{Clean extraction.}
Call an unmarked site clean. Two clean nonblank labels at the
same \(u,t\) agree, even at different transverse positions.
Indeed each has four supporters, so their supporter sets intersect
in at least three indices. Each clean root has at most one dirty
provider; after excluding both dirty indices a common wholly clean
supporter remains. Its transverse equality checks move one label
one unit toward the other at the preceding physical time.
Iterating decreases their transverse \(\ell^1\) distance until
they occupy the same site, proving equality of the full labels.

Augment any TI stationary noisy law by its two-sided iid innovations,
using consistent finite time windows, and append \(E^*\).
Let \(y_t(u)\) be the unique nonblank label found at a clean
transverse position, or \(0\) if none exists. A clean live witness
has a wholly clean successful provider, which supplies its same
old center and every positive parent prescribed by its witness.
Thus \(y_t\le y_{t-1}\) in deletion order. Time stationarity makes
their nonblank probabilities equal, so they agree almost surely.
Countability makes the extracted field time constant and
\(H_{0,e}\)-fixed, with a fully \(\Z^2\)-TI law.
On totality, source clause (i) makes it blank. Every nonblank
vertex in every stationary TI noisy margin is therefore marked.

Now take an arbitrary TI stationary canonical common-noise pair
and augment both replicas by the same innovations and marks.
A discrepancy is unflagged, since shared replacement gives equal
letters. Distinct deterministic outputs require a discrepant
parent. Each discrepancy has a nonblank margin, hence is marked.
Tracing parents into the two-sided past would give an infinite
marked nonflagged path, excluded by~\eqref{eq:marked-path}.
The pair is diagonal. This is the arbitrary-pair quantifier in
the usual stationary uniqueness argument made explicit.

\emph{Separation on nontotality.}
Source clause (ii) gives a positive-activity \(H_{0,e}\)-fixed
reference law and also the blank reference law. Clone either
reference constantly in \(P\). It is fixed by \(G_e\).
Start the noisy field at its reference and the mark field in its
stationary least law, independently of the reference. The closed
relation
\[
 X_t(v)\ne\operatorname{Ref}(v)\ \Longrightarrow\ E_t(v)=1
\]
is preserved: at a clean site four wholly clean providers have
the correct reference inputs, including when the reference is
blank. Fixed-marginal Ces\`aro limits give stationary TI noisy
lifts with this error domination.

For fixed \(u\), the transverse mark field is mixing: approximate
its indicator in \(L^1\) by a finite-past function of finitely many
flags, whose distant translates are independent. Its transverse
cube averages have variance tending to zero and mean \(p_*<1/2\).
Let \(\delta_R\) be the probability that at least half the sites
of the radius-\(R\) transverse cube are marked; then
\(\delta_R\to0\). Strict majority decodes \(\operatorname{Ref}(u)\)
except on this event. If the active reference has nonblank
probability \(a_*>0\), the two noisy lifts have local decoded
nonblank probabilities at least \(a_*-\delta_R\) and at most
\(\delta_R\). Some finite \(R\) separates them. This observable
is needed only to prove distinctness; the compiler uses the
finite source table.
\end{proof}

\subsection{The simultaneous count reduction}
One reduction family establishes hardness for every finite upper
and lower threshold.
\begin{maintheorem}[Phase-count completeness under \(C_T\)]\label{thm:completeness}
Assume \(C_T\). At \(D=7\) and the single rate
\(\eps_1=160000^{-101}\), in each input category separately,
\(\{M:N(M)\le r\}\) is \(\Pi^0_2\)-complete for every finite
\(r\ge1\), and \(\{M:N(M)\ge k\}\) is \(\Sigma^0_2\)-complete
for every finite \(k\ge2\), under computable many-one reductions.
One computable family has \(N=1\) on total programs and
\(N=\frakc\) on nontotal programs.
\end{maintheorem}
\begin{proof}
Run \(G_e\) independently in each seventh-coordinate slice,
with the same rate \(\eps_1\), and call the result \(C_e\).
On totality, a TI stationary common-noise pair restricts on each
slice to such a pair for \(G_e\), hence to equal replicas by
Lemma~\ref{lem:totality-controller}. Countably many slices make
the full pair diagonal. Given any two TI stationary laws, evolve
their product using common noise and take a Ces\`aro limit.
The marginals stay fixed and the limit is a stationary TI pair.
Diagonality equates its marginals. Existence follows from
Lemma~\ref{lem:compact}, so \(N(C_e)=1\).

On nontotality choose two distinct stationary TI slice laws
\(\mu_0,\mu_1\), and a bounded real local \(h\) separating
their means. For \(0\le p\le1\), put
\[
 \mu_p=(1-p)\mu_0+p\mu_1,\qquad
 \nu_p=\bigotimes_{j\in\Z}\mu_p.
\]
These laws are stationary and fully TI. Translates of any local
observable far enough apart along \(e_7\) depend on disjoint
sets of independent slices. Their Ces\`aro average has variance
tending to zero. A decomposition of \(\nu_p\) into two TI laws
with distinct local means would give a positive lower bound on
that variance. Thus \(\nu_p\) is extreme among all TI laws.
The \(h\)-means distinguish continuum many \(p\), and
Lemma~\ref{lem:compact} bounds their number by \(\frakc\).
Hence
\begin{equation}\label{eq:one-continuum}
 \varphi_e\text{ total}\ \Longrightarrow\ N(C_e)=1,\qquad
 \varphi_e\text{ nontotal}\ \Longrightarrow\ N(C_e)=\frakc.
\end{equation}
The classical range stays at most \(11\); the finite alphabet
may depend on \(e\).

The set \(\mathrm{TOT}=\{e:\forall n\ \exists s\
[\varphi_e(n)\text{ halts by }s]\}\) is \(\Pi^0_2\)-complete.
For a decidable relation \(R(a,n,t)\), compile a program on
input \(n\) that searches for a \(t\) satisfying \(R(a,n,t)\).
It is total exactly when \(\forall n\,\exists t\,R(a,n,t)\).
Complementation gives \(\Sigma^0_2\)-completeness of
\(\overline{\mathrm{TOT}}\). Equation~\eqref{eq:one-continuum}
reduces \(\mathrm{TOT}\) to every finite upper cut \(N\le r\),
and its complement to every lower cut \(N\ge k\), \(k\ge2\).
The unconditional upper bounds of Theorem~\ref{thm:counts}
give the matching membership statements.

For quantum inputs, implement \(C_e\) by the
dephase/copy/lookup/prepare circuit of
Section~\ref{sec:preservation}, followed by whole-site
depolarization at \(\eps_1\). Every finite output block is
basis diagonal after one round. Thus every stationary quantum
state is diagonal, and its stationary TI set is affinely exactly
the classical one. Equation~\eqref{eq:one-continuum} holds for
the entire TI stationary quantum set. Theorem~\ref{thm:counts} also
gives quantum membership. On the quantum presentation promise,
membership means a total arithmetic formula agrees on valid
inputs; every reduction here outputs a valid exact circuit.
\end{proof}

\subsection{Stationary preservation from diagonal couplings}
For the general stationary Rice clause, discard the \(E\) coordinate:
\[
 \sigma_{\mathrm{stat}}=(N,N_{\ord},N_{\Br},L),\qquad
 \sigma_{\mathrm{stat},\infty}=(\frakc,\frakc,\frakc,\mathrm{true}).
\]

\begin{proposition}\label{prop:stationary-diagonal}
At any common rate, coupling a host to a deterministic TI
controller followed by uniform replacement and having the
diagonal-coupling property, using one whole-pair replacement,
preserves \(\sigma_{\mathrm{stat}}\), in both categories.
Host projection is an affine homeomorphism of TI stationary sets
and preserves \(\ell\) and every \(b_q\) per state.
\end{proposition}
\begin{proof}
The controller has a unique TI stationary law \(\varpi\), by the
fixed-marginal Ces\`aro argument above. Start its common-noise
pair from \((\varpi,\mathbf0)\), and let \(d_m\) be the one-site
disagreement after \(m\) rounds. Every subsequential Ces\`aro
pair limit is TI stationary and diagonal. Since disagreement
is a cylinder event,
\[
 \frac1T\sum_{m<T}d_m\longrightarrow0.
\]
Otherwise a compactness limit would have positive disagreement.
Choose an unbounded subsequence \(m_j\) with \(d_{m_j}\to0\).

Every TI stationary host state \(\mu\) has a stationary joint
lift. In any such lift \(\zeta\), the controller marginal is
\(\varpi\). Retain its host start and future innovations, but
reset the comparison controller to \(\mathbf0\).
The controller pair has exactly the law started from
\((\varpi,\mathbf0)\), regardless of initial host correlations.
For the finite-past comparison \(W_m(\mu)\) this gives
\begin{equation}\label{eq:stationary-reset}
 d_O(\zeta,W_m(\mu))\le |O_C|d_m.
\end{equation}
Along \(m_j\), all stationary lifts of \(\mu\) agree.
The continuous affine onto projection is therefore an affine
homeomorphism and bijects the full extreme sets.

Quantumly, a stationary joint state is diagonal on the
measure-and-prepare controller. On a finite input cone write it
as
\(\sum_c\rho_H^c\otimes|c\rangle\langle c|\), with positive
subnormalized host matrices and traces given by the
\(\varpi\)-marginal. At fixed innovations the channel factors
between tracks. The actual and reset outputs have the same host
images of \(\rho_H^c\); their normalized trace distance is at
most the total weight of strings whose controller outputs differ
on \(O_C\). Averaging gives exactly the classical mismatch
bound \(|O_C|d_m\). This proves~\eqref{eq:stationary-reset}
also with initially entangled host data.

For joint observables \(f,g\), let \(r\) bound the number of
controller sites in either support. The covariance error in
\eqref{eq:covariance-error} is now at most
\(8\|f\|\|g\|r d_m\). At fixed \(m\), distant finite cones
factor through local host observables \(G_m^f,G_m^g\);
quantumly the factorization follows from the independent
product environments. If the host state lacks \(\ell\), send
separation to infinity and then \(m=m_j\to\infty\) to exclude
it in the lift. Host observables give the converse.
For each \(q\ne0\), expand the weighted second moment:
only \(O_m(|B_n|)\) pairs have intersecting cones, and the
remaining products are host products. The stationary/reset
error is bounded uniformly in \(n\) by \(8\|f\|^2r d_m\).
First \(n\to\infty\), then \(m=m_j\to\infty\), preserves
each \(b_q\), exactly as in Proposition~\ref{prop:preservation}.
Thus \(N_{\ord},N_{\Br}\) and \(L\) are preserved as well.
\end{proof}

This proposition controls stationary controller marginals; its
conclusion concerns \(\sigma_{\mathrm{stat}}\), excluding
all-law ergodicity. A constant bath sharing the replacement
flags is a trivial diagonal-coupling controller: replicas agree
after one step. Its invariance follows from this proposition.
The product lift \(\mu\otimes\omega_{\mathrm B}\) in
Proposition~\ref{prop:bath} refers to a bath with independent
noise in the comparison universe.

\subsection{The stationary Rice degrees}
\begin{theorem}[Stationary Rice degrees under \(C_T\)]\label{thm:stationary-degrees}
Assume \(C_T\), \(D=7\) and rate \(\eps_1\). In each input
category, every nontrivial predicate \(P\) of
\(\sigma_{\mathrm{stat}}\) is \(\Pi^0_2\)-hard if it is false
at \(\sigma_{\mathrm{stat},\infty}\), and \(\Sigma^0_2\)-hard
if it is true there. There is no nontrivial
\(\Delta^0_2\) predicate of this datum. If such a predicate
separately belongs to \(\Pi^0_2\), or to \(\Sigma^0_2\), it is
complete for that class.
\end{theorem}
\begin{proof}
Pair two copies of \(G_e\), slice in the seventh direction,
and apply the parity conjugacy~\eqref{eq:gauge}, all at
\(\eps_1\). Denote the resulting controller by \(Q_e\).
On totality the diagonal-coupling property survives slicing
and pairing, by restriction to slices and to corresponding
registers of the two replicas. For parity conjugacy, pull back
a fully TI stationary pair through \(T\) in both replicas.
The pulled-back law is stationary, invariant under even
translations, and each odd translation acts on it by the global
swap in both replicas. Averaging it with its swap image makes
it fully TI, hence diagonal by the ungauged property.
Each component of this positive average is diagonal, so the
original pair is diagonal. Proposition~\ref{prop:stationary-diagonal}
therefore preserves every host's \(\sigma_{\mathrm{stat}}\).

On nontotality, let \(h\) be any real local separator of the
two stationary slice laws, with mean gap \(\delta>0\).
Fixed-marginal Ces\`aro averaging gives a paired joining
\(\xi_{10}\); let \(\xi_{01}\) be its swap. Set
\[
 \xi_p=(1-p)\xi_{10}+p\xi_{01},\quad
 g=h(\text{first})-h(\text{second}),\quad
 \delta_p=(1-2p)\delta,\quad 0<p<\tfrac12.
\]
Use iid \(\xi_p\) slices and average the two parity-conjugate
components as in Proposition~\ref{prop:crystal}. For every local
observable the component cube-average variance is \(O(1/n)\),
because only \(O(|B_n|^2/n)\) pairs share a slice. Their parity
means average to the same limit. Mixture variance therefore
proves that the TI average \(\rho_p\) is extreme among all TI
laws, and among all TI quantum states in its diagonal realization.

Here \(g\) may read several sites inside one slice. Use the
local observable \(f=T^*g\), anchored at the origin. Swap
oddness gives
\[
 T^*\tau_x(T^*g)=(-1)^{\sum_i x_i}\tau_xg.
\]
The staggered component means are \(\delta_p\) and
\(-\delta_p\), so the \(q_\pi\) Bragg moment is at least
\(\delta_p^2\). Moreover
\(\rho_p(f\tau_{2e_7}f)=\delta_p^2\); independence of the two
whole slices proves this identity and distinguishes continuum
many \(p\). Thus \(Q_e\) has continuum many Bragg-ordered
extreme TI stationary states.

Write \(H\boxdot Q_e\) for the product of the host and controller
updates followed by one whole-pair replacement.
For every host \(H\), projection of its stationary set onto
the controller is onto.
The fiber over each extreme \(\rho_p\) is a compact face and
has an extreme joint lift by Lemma~\ref{lem:compact}.
The lift retains the controller Bragg moment. By
Lemma~\ref{lem:two-order} and the continuum upper bound, every
nontotal product has datum \(\sigma_{\mathrm{stat},\infty}\).
Total products have datum \(\sigma_{\mathrm{stat}}(H)\).
The quantum assertions follow from the same diagonal controller
and the trace-distance argument just proved.

A known nontotal program realizes
\(\sigma_{\mathrm{stat},\infty}\). Nontriviality supplies a
finite valid host \(H\) with the opposite verdict; fix it.
The computable map \(e\mapsto H\boxdot Q_e\) reduces
\(\mathrm{TOT}\) to \(P\) when its universal verdict is false,
and \(\overline{\mathrm{TOT}}\) to \(P\) when it is true.
Only finite tables and circuits enter this map.

For completeness, hierarchy strictness forces the matching
orientations. The relation
\[
 V(e,x)\iff\forall n\ \exists s\
 [\text{program }e\text{ on }(x,n)\text{ halts by }s]
\]
is universal for \(\Pi^0_2\). Its diagonal complement
\(\neg V(x,x)\) is \(\Sigma^0_2\) and cannot be \(\Pi^0_2\):
an index \(d\) for such a definition would give
\(\neg V(d,d)=V(d,d)\). Hence \(\mathrm{TOT}\notin\Sigma^0_2\)
and \(\overline{\mathrm{TOT}}\notin\Pi^0_2\).
A \(\Pi^0_2\) predicate with true universal verdict would
place \(\overline{\mathrm{TOT}}\) in \(\Pi^0_2\), a contradiction.
It must have false verdict and is \(\Pi^0_2\)-complete.
The dual statement holds for \(\Sigma^0_2\); belonging to both
classes is impossible for a nontrivial predicate.
On the quantum promise domain these preimages are ordinary
arithmetic predicates because every output presentation is valid.
\end{proof}

The concrete complete predicates established here are the total-count
cuts of Theorem~\ref{thm:completeness}. For \(N_{\ord},N_{\Br}\)
and \(L\), the general clause supplies one-sided hardness and the
absence of nontrivial \(\Delta^0_2\) members; an arithmetic upper
bound would be needed to invoke its matching-completeness clause.
Exact \(N=k\), \(k\ge2\), is \(\Pi^0_2\)-hard if realized and
has the unconditional \(D_2(\Sigma^0_2)\) upper bound, while
realizability remains open. Finite and infinite total count retain
their \(\Sigma^0_3\) and \(\Pi^0_3\) upper bounds.
